% Options for packages loaded elsewhere
\PassOptionsToPackage{unicode}{hyperref}
\PassOptionsToPackage{hyphens}{url}
\PassOptionsToPackage{dvipsnames,svgnames,x11names}{xcolor}
\documentclass[
  spanish,
  10pt,
  a4paper,
]{article}
\usepackage{xcolor}
\usepackage[top=2.2cm,bottom=2.2cm,left=2.1cm,right=2.1cm]{geometry}
\usepackage{amsmath,amssymb}
\setcounter{secnumdepth}{5}
\usepackage{iftex}
\ifPDFTeX
  \usepackage[T1]{fontenc}
  \usepackage[utf8]{inputenc}
  \usepackage{textcomp} % provide euro and other symbols
\else % if luatex or xetex
  \usepackage{unicode-math} % this also loads fontspec
  \defaultfontfeatures{Scale=MatchLowercase}
  \defaultfontfeatures[\rmfamily]{Ligatures=TeX,Scale=1}
\fi
\usepackage{lmodern}
\ifPDFTeX\else
  % xetex/luatex font selection
\fi
% Use upquote if available, for straight quotes in verbatim environments
\IfFileExists{upquote.sty}{\usepackage{upquote}}{}
\IfFileExists{microtype.sty}{% use microtype if available
  \usepackage[]{microtype}
  \UseMicrotypeSet[protrusion]{basicmath} % disable protrusion for tt fonts
}{}
\makeatletter
\@ifundefined{KOMAClassName}{% if non-KOMA class
  \IfFileExists{parskip.sty}{%
    \usepackage{parskip}
  }{% else
    \setlength{\parindent}{0pt}
    \setlength{\parskip}{6pt plus 2pt minus 1pt}}
}{% if KOMA class
  \KOMAoptions{parskip=half}}
\makeatother
\usepackage{color}
\usepackage{fancyvrb}
\newcommand{\VerbBar}{|}
\newcommand{\VERB}{\Verb[commandchars=\\\{\}]}
\DefineVerbatimEnvironment{Highlighting}{Verbatim}{commandchars=\\\{\}}
% Add ',fontsize=\small' for more characters per line
\newenvironment{Shaded}{}{}
\newcommand{\AlertTok}[1]{\textcolor[rgb]{1.00,0.00,0.00}{\textbf{#1}}}
\newcommand{\AnnotationTok}[1]{\textcolor[rgb]{0.38,0.63,0.69}{\textbf{\textit{#1}}}}
\newcommand{\AttributeTok}[1]{\textcolor[rgb]{0.49,0.56,0.16}{#1}}
\newcommand{\BaseNTok}[1]{\textcolor[rgb]{0.25,0.63,0.44}{#1}}
\newcommand{\BuiltInTok}[1]{\textcolor[rgb]{0.00,0.50,0.00}{#1}}
\newcommand{\CharTok}[1]{\textcolor[rgb]{0.25,0.44,0.63}{#1}}
\newcommand{\CommentTok}[1]{\textcolor[rgb]{0.38,0.63,0.69}{\textit{#1}}}
\newcommand{\CommentVarTok}[1]{\textcolor[rgb]{0.38,0.63,0.69}{\textbf{\textit{#1}}}}
\newcommand{\ConstantTok}[1]{\textcolor[rgb]{0.53,0.00,0.00}{#1}}
\newcommand{\ControlFlowTok}[1]{\textcolor[rgb]{0.00,0.44,0.13}{\textbf{#1}}}
\newcommand{\DataTypeTok}[1]{\textcolor[rgb]{0.56,0.13,0.00}{#1}}
\newcommand{\DecValTok}[1]{\textcolor[rgb]{0.25,0.63,0.44}{#1}}
\newcommand{\DocumentationTok}[1]{\textcolor[rgb]{0.73,0.13,0.13}{\textit{#1}}}
\newcommand{\ErrorTok}[1]{\textcolor[rgb]{1.00,0.00,0.00}{\textbf{#1}}}
\newcommand{\ExtensionTok}[1]{#1}
\newcommand{\FloatTok}[1]{\textcolor[rgb]{0.25,0.63,0.44}{#1}}
\newcommand{\FunctionTok}[1]{\textcolor[rgb]{0.02,0.16,0.49}{#1}}
\newcommand{\ImportTok}[1]{\textcolor[rgb]{0.00,0.50,0.00}{\textbf{#1}}}
\newcommand{\InformationTok}[1]{\textcolor[rgb]{0.38,0.63,0.69}{\textbf{\textit{#1}}}}
\newcommand{\KeywordTok}[1]{\textcolor[rgb]{0.00,0.44,0.13}{\textbf{#1}}}
\newcommand{\NormalTok}[1]{#1}
\newcommand{\OperatorTok}[1]{\textcolor[rgb]{0.40,0.40,0.40}{#1}}
\newcommand{\OtherTok}[1]{\textcolor[rgb]{0.00,0.44,0.13}{#1}}
\newcommand{\PreprocessorTok}[1]{\textcolor[rgb]{0.74,0.48,0.00}{#1}}
\newcommand{\RegionMarkerTok}[1]{#1}
\newcommand{\SpecialCharTok}[1]{\textcolor[rgb]{0.25,0.44,0.63}{#1}}
\newcommand{\SpecialStringTok}[1]{\textcolor[rgb]{0.73,0.40,0.53}{#1}}
\newcommand{\StringTok}[1]{\textcolor[rgb]{0.25,0.44,0.63}{#1}}
\newcommand{\VariableTok}[1]{\textcolor[rgb]{0.10,0.09,0.49}{#1}}
\newcommand{\VerbatimStringTok}[1]{\textcolor[rgb]{0.25,0.44,0.63}{#1}}
\newcommand{\WarningTok}[1]{\textcolor[rgb]{0.38,0.63,0.69}{\textbf{\textit{#1}}}}
\usepackage{longtable,booktabs,array}
\newcounter{none} % for unnumbered tables
\usepackage{calc} % for calculating minipage widths
% Correct order of tables after \paragraph or \subparagraph
\usepackage{etoolbox}
\makeatletter
\patchcmd\longtable{\par}{\if@noskipsec\mbox{}\fi\par}{}{}
\makeatother
% Allow footnotes in longtable head/foot
\IfFileExists{footnotehyper.sty}{\usepackage{footnotehyper}}{\usepackage{footnote}}
\makesavenoteenv{longtable}
\ifLuaTeX
\usepackage[bidi=basic,shorthands=off]{babel}
\else
\usepackage[bidi=default,shorthands=off]{babel}
\fi
\ifLuaTeX
  \usepackage{selnolig} % disable illegal ligatures
\fi
\setlength{\emergencystretch}{3em} % prevent overfull lines
\providecommand{\tightlist}{%
  \setlength{\itemsep}{0pt}\setlength{\parskip}{0pt}}
% Shared visual layer for the generated Kaleido FlowTwin reports.
\usepackage{helvet}
\renewcommand{\familydefault}{\sfdefault}
\usepackage{microtype}
\usepackage{xcolor}
\usepackage{titlesec}
\usepackage{fancyhdr}
\usepackage{enumitem}
\usepackage{booktabs}
\usepackage{xurl}
\usepackage{lastpage}
\usepackage{etoolbox}

\definecolor{KaleidoNavy}{HTML}{0A2342}
\definecolor{KaleidoBlue}{HTML}{176B87}
\definecolor{KaleidoCyan}{HTML}{3BB7C4}
\definecolor{KaleidoMint}{HTML}{DDF5F1}
\definecolor{KaleidoFog}{HTML}{F2F5F7}
\definecolor{KaleidoSlate}{HTML}{5D6A7A}
\definecolor{KaleidoCoral}{HTML}{F26B5B}

\titleformat{\section}
  {\Large\bfseries\color{KaleidoNavy}}
  {\thesection}{0.75em}{}[\vspace{0.2em}\color{KaleidoCyan}\titlerule]
\titleformat{\subsection}
  {\large\bfseries\color{KaleidoBlue}}
  {\thesubsection}{0.65em}{}
\titleformat{\subsubsection}
  {\normalsize\bfseries\color{KaleidoNavy}}
  {\thesubsubsection}{0.55em}{}
\titlespacing*{\section}{0pt}{2.2em}{1em}
\titlespacing*{\subsection}{0pt}{1.6em}{0.65em}

\pagestyle{fancy}
\fancyhf{}
\fancyhead[L]{\small\bfseries\color{KaleidoNavy}KALEIDO FLOWTWIN}
\fancyhead[R]{\small\color{KaleidoSlate}DOSSIER TECNICO}
\fancyfoot[L]{\scriptsize\color{KaleidoSlate}EVOCON Solutions GmbH}
\fancyfoot[R]{\scriptsize\color{KaleidoSlate}\thepage/\pageref*{LastPage}}
\renewcommand{\headrulewidth}{0.6pt}
\renewcommand{\headrule}{\hbox to\headwidth{\color{KaleidoCyan}\leaders\hrule height \headrulewidth\hfill}}
\renewcommand{\footrulewidth}{0pt}
\setlength{\headheight}{20pt}

\setlist[itemize]{leftmargin=1.4em,itemsep=0.25em,topsep=0.35em}
\setlist[enumerate]{leftmargin=1.7em,itemsep=0.25em,topsep=0.35em}
\setlength{\parindent}{0pt}
\setlength{\parskip}{0.55em}
\setlength{\emergencystretch}{3em}
\renewcommand{\arraystretch}{1.18}

\AtBeginEnvironment{quote}{%
  \color{KaleidoNavy}\itshape
  \setlength{\leftskip}{1em}
  \setlength{\rightskip}{1em}
}

\makeatletter
\renewcommand{\maketitle}{%
  \hypersetup{pageanchor=false}
  \begin{titlepage}
    \pagecolor{KaleidoNavy}
    \color{white}
    \vspace*{1.4cm}
    {\color{KaleidoCyan}\rule{2.4cm}{2.5pt}\par}
    \vspace{1.1cm}
    {\Huge\bfseries\@title\par}
    \vspace{0.65cm}
    {\Large\color{KaleidoMint}Kaleido FlowTwin\par}
    \vfill
    {\large Álvaro Schwiedop Souto\par}
    \vspace{0.15cm}
    {\normalsize\color{KaleidoMint}Industrial AI \& Predictive Quality Lead\par}
    \vspace{0.15cm}
    {\normalsize\color{KaleidoMint}EVOCON Solutions GmbH\par}
    \vspace{0.25cm}
    {\normalsize\color{KaleidoMint}\@date\par}
  \end{titlepage}
  \hypersetup{pageanchor=true}
  \nopagecolor
  \color{black}
}
\makeatother
\usepackage{bookmark}
\IfFileExists{xurl.sty}{\usepackage{xurl}}{} % add URL line breaks if available
\urlstyle{same}
% fallback for those not using the hyperref driver hyperxmp:
\makeatletter
\@ifundefined{xmpquote}{\newcommand{\xmpquote}[1]{#1}}{}
\makeatother
\hypersetup{
  pdftitle={Plan del MVP Kaleido FlowTwin},
  pdflang={es},
  colorlinks=true,
  linkcolor={KaleidoBlue},
  filecolor={Maroon},
  citecolor={Blue},
  urlcolor={KaleidoBlue},
  pdfcreator={LaTeX via pandoc}}

\title{Plan del MVP Kaleido FlowTwin}
\author{}
\date{15 julio 2026}

\begin{document}
\maketitle

{
\setcounter{tocdepth}{3}
\tableofcontents
}
\section{Plan del MVP Kaleido
FlowTwin}\label{plan-del-mvp-kaleido-flowtwin}

\subsection{Objetivo y decision}\label{objetivo-y-decision}

Construir un piloto offline/read-only para predecir tiempo restante y
riesgo de desviacion en una operacion/turno de Trace Port. El plan
admite que el historico sea insuficiente: en ese caso termina con
instrumentacion, process intelligence y simulacion, sin inventar un
predictor.

\subsection{Estados de salida}\label{estados-de-salida}

{\def\LTcaptype{none} % do not increment counter
\begin{longtable}[]{@{}lll@{}}
\toprule\noalign{}
Estado & Significado & Siguiente decision \\
\midrule\noalign{}
\endhead
\bottomrule\noalign{}
\endlastfoot
\texttt{GO\_PREDICTIVE} & hay volumen, calidad, outcomes y senal &
piloto shadow 4-8 semanas \\
\texttt{INSTRUMENT\_FIRST} & proceso util pero dato insuficiente &
activar captura minima y reevaluar \\
\texttt{BASELINE\_ONLY} & modelo simple gana & productizar baseline, no
JEPA \\
\texttt{NO\_SIGNAL} & no hay mejora ni KPI accionable & parar o cambiar
caso \\
\end{longtable}
}

Todos son resultados validos del piloto.

\subsection{Fase 0 - Alineamiento
(reunion)}\label{fase-0---alineamiento-reunion}

Duracion: 60-90 minutos.

Entregables:

\begin{itemize}
\tightlist
\item
  proceso candidato;
\item
  comprador/usuario;
\item
  unidad de decision;
\item
  desviacion material;
\item
  acciones permitidas;
\item
  KPI y coste aproximado;
\item
  propietario de datos;
\item
  export minimo acordado.
\end{itemize}

Gate: una pregunta predictiva con timestamp de decision y outcome
verificable.

\subsection{Fase 1 - Data Evidence
Scan}\label{fase-1---data-evidence-scan}

Duracion: 3-5 dias tras recibir muestra.

Entrada: 3-5 operaciones anonimizadas y diccionario/export.

Tareas:

\begin{itemize}
\tightlist
\item
  schema y relaciones;
\item
  timezone, IDs, duplicados y revisiones;
\item
  action/context/outcome audit;
\item
  cobertura de plan, actual y outcome;
\item
  inventario de eventos/modalidades;
\item
  riesgo GDPR/seguridad;
\item
  estimacion de operaciones utilizables y diversidad.
\end{itemize}

Salida:

\begin{itemize}
\tightlist
\item
  \texttt{data\_readiness\_report.html/pdf};
\item
  canonical mapping;
\item
  leakage report;
\item
  decision \texttt{go/instrument/no-fit};
\item
  presupuesto de piloto.
\end{itemize}

Gate: se puede reconstruir al menos una operacion sin mirar el futuro.

\subsection{Fase 2 - Proceso y
baseline}\label{fase-2---proceso-y-baseline}

Duracion: semana 1-2.

Tareas:

\begin{itemize}
\tightlist
\item
  object-centric event log;
\item
  process map, variantes, esperas y conformance;
\item
  plan vs actual;
\item
  naive/median/survival/boosting;
\item
  split cronologico y por proyecto;
\item
  metricas operativas preliminares.
\end{itemize}

Gate:

\begin{itemize}
\tightlist
\item
  artefactos reproducibles;
\item
  umbrales solo de validation;
\item
  baseline aporta informacion por encima de plan/mediana;
\item
  explicacion revisable por un operador.
\end{itemize}

\subsection{Fase 3 - Modelo
secuencial}\label{fase-3---modelo-secuencial}

Duracion: semana 2-4, si Fase 2 pasa.

Tareas:

\begin{itemize}
\tightlist
\item
  prefijos causales a varios prediction points;
\item
  GRU/TCN/ProcessTransformer;
\item
  object-centric graph baseline;
\item
  quantiles/calibracion;
\item
  error por grupo y drift temporal.
\end{itemize}

Gate: mejora material sobre boosting en al menos una metrica operativa
sin empeorar de forma inaceptable calibracion o peor grupo.

\subsection{Fase 4 - Event-JEPA
experimental}\label{fase-4---event-jepa-experimental}

Duracion: semana 4-5, condicionada al volumen.

Tareas:

\begin{itemize}
\tightlist
\item
  future-embedding prediction 2/4/8 h;
\item
  target encoder/anticollapse;
\item
  actions vs context;
\item
  ablations correct/shuffled/no-action;
\item
  frozen probes y fine-tuning;
\item
  tres seeds.
\end{itemize}

Gate de promocion:

\begin{itemize}
\tightlist
\item
  delta positivo frente a mejor secuencial;
\item
  lead time o falsas alertas mejores;
\item
  mejora sostenida por seed/holdout;
\item
  correct actions \textgreater{} shuffled para action claim;
\item
  latencia y calibracion aceptables.
\end{itemize}

Si falla, se elimina del producto.

\subsection{Fase 5 - Escenarios}\label{fase-5---escenarios}

Duracion: semana 4-6, en paralelo conceptual con Fase 4.

Ruta A, prioritaria:

\begin{itemize}
\tightlist
\item
  simulacion discreta de tareas, recursos, colas, pausas y calendario;
\item
  calibracion con tiempos observados;
\item
  escenarios definidos con experto;
\item
  sensibilidad y rangos, no una cifra unica.
\end{itemize}

Ruta B, solo con acciones reales:

\begin{itemize}
\tightlist
\item
  ranking de acciones observadas;
\item
  soporte/uncertainty guard;
\item
  offline policy evaluation conservadora;
\item
  advisory-only.
\end{itemize}

Gate: escenarios plausibles, dentro de restricciones y con hipotesis
visibles.

\subsection{Fase 6 - Shadow MVP}\label{fase-6---shadow-mvp}

Duracion: semana 5-6 para replay; 4-8 semanas prospectivas si Kaleido
continua.

Entregables:

\begin{itemize}
\tightlist
\item
  dashboard de operacion;
\item
  P50/P90 y riesgo 2/4/8 h;
\item
  timeline y fuente de cada alerta;
\item
  motivos/objetos afectados;
\item
  abstencion;
\item
  export PDF/CSV/API;
\item
  model/data cards;
\item
  informe tecnico y ejecutivo.
\end{itemize}

Gate de piloto prospectivo, a acordar:

\begin{itemize}
\tightlist
\item
  lead time minimo;
\item
  limite de falsas alertas/turno;
\item
  cobertura P90;
\item
  mejora vs plan/baseline;
\item
  utilidad operatoria;
\item
  disponibilidad/latencia.
\end{itemize}

\subsection{Cronograma comercial
propuesto}\label{cronograma-comercial-propuesto}

\begin{Shaded}
\begin{Highlighting}[]
\NormalTok{Reunion        seleccionar proceso y muestra}
\NormalTok{Dias 1{-}5       Data Evidence Scan}
\NormalTok{Semana 2       mapa de proceso + baseline}
\NormalTok{Semana 3{-}4     predictor y calibracion}
\NormalTok{Semana 5       escenarios + dashboard}
\NormalTok{Semana 6       replay, informe y go/no{-}go}
\NormalTok{Posterior      shadow prospectivo y productizacion}
\end{Highlighting}
\end{Shaded}

No vender \texttt{4-6\ semanas} como garantia de precision. Es plazo
para obtener una decision basada en evidencia.

\subsection{Matriz de exito}\label{matriz-de-exito}

\subsubsection{Tecnico}\label{tecnico}

\begin{itemize}
\tightlist
\item
  integridad temporal;
\item
  holdout sin leakage;
\item
  intervalo calibrado;
\item
  mejora frente a baseline;
\item
  robustez por proyecto/carga;
\item
  latencia y trazabilidad.
\end{itemize}

\subsubsection{Operativo}\label{operativo}

\begin{itemize}
\tightlist
\item
  alertas antes de que el operador ya conozca el problema;
\item
  volumen de falsas alarmas tolerable;
\item
  razon de alerta comprensible;
\item
  accion disponible en ese momento;
\item
  dashboard no duplica trabajo.
\end{itemize}

\subsubsection{Comercial}\label{comercial}

\begin{itemize}
\tightlist
\item
  comprador identificado;
\item
  coste del problema medible;
\item
  integracion con producto Kaleido;
\item
  packaging y precio potencial;
\item
  IP/datos/soporte acordables;
\item
  ruta a mas clientes sin mezclar datos.
\end{itemize}

\subsection{Riesgos principales}\label{riesgos-principales}

{\def\LTcaptype{none} % do not increment counter
\begin{longtable}[]{@{}ll@{}}
\toprule\noalign{}
Riesgo & Mitigacion \\
\midrule\noalign{}
\endhead
\bottomrule\noalign{}
\endlastfoot
historico pobre & instrument-first, public/synthetic solo para
plumbing \\
pocos eventos negativos & remaining time, survival, top-risk; no
clasificador forzado \\
plan sobreescrito & revisionado y cutoff temporal \\
IDs no enlazables & resolver mapping antes de modelar \\
atajos por cliente/proyecto & holdout, ablations y worst-group \\
drift de proceso & frozen reference + shadow, reset por redisenos \\
accion confundida con contexto & contrato y shuffled-action test \\
ROI no disponible & separar metricas tecnicas de simulacion economica \\
duplicar Trace Port & integrar la prediccion dentro del producto \\
scope creep 3D/robotica & backlog separado, no primera version \\
\end{longtable}
}

\subsection{Backlog posterior}\label{backlog-posterior}

\begin{itemize}
\tightlist
\item
  Shipping Board remaining-time/exception transfer;
\item
  Freight Intelligence ETA probabilistica;
\item
  fotos de incidencia con DINO frozen features;
\item
  root-cause cross-product;
\item
  optimizacion de recursos bajo restricciones;
\item
  LiDAR/vision para ocupacion y maniobras;
\item
  aprendizaje federado o por cliente con adapters.
\end{itemize}

\end{document}
